\section{Unitary Evaporation and Page-Like Behavior}

The collision model gives a unitary purification of the emitted radiation process. At each step a fresh bin begins in $|0\rangle$, interacts with the core by $U=\exp(-iH_{\rm coll}\Delta t)$, and then becomes part of the outgoing radiation register. Iterating this procedure produces a pure state on
\begin{equation}
  \Hcore\otimes \bigotimes_{j=1}^{n}\mathcal{H}_{{\rm rad},j}.
\end{equation}

In the explicit simulations we use the reduced core channel generated by the same unitary. This avoids storing the full radiation history, whose dimension grows as $(M+1)^n$. The Renyi-2 radiation entropy is still computed exactly, because for the global pure state
\begin{equation}
  S^{(2)}_{\rm rad}(n)
  =
  S^{(2)}_{\rm core}(n)
  =
  -\log {\rm Tr}\,\rho_{\rm core}(n)^2.
\end{equation}

\subsection{One-Channel Failure}

A naive one-channel Hamiltonian does not work well. The reason is structural. If a transition operator maps
\begin{equation}
  X_m:\mathcal{H}_m\rightarrow\mathcal{H}_{m+1},
\end{equation}
then its rank is at most $D_{m+1}$. In an evaporating model $D_{m+1}<D_m$, so much of $\mathcal{H}_m$ is dark to the emission operator. The bright component emits, while the dark component remains trapped. Numerically this produces front-loaded emission and no robust acceleration.

This failure is useful. It shows that negative heat capacity alone is not enough. An evaporator also needs enough outgoing channel capacity, or enough internal scrambling, to avoid dark-subspace trapping. In the present model we implement this by allowing several emitted labels per time bin.

\subsection{Step 2 Result}

The polished Step 2 comparison uses eight emitted labels, weak coupling, and twelve random seeds. The convex run uses a strongly convex entropy profile; the control uses a linear entropy profile. Both use the same collision architecture.

For the convex case,
\begin{equation}
  {\rm curvature}=3,\qquad M=8,\qquad g=0.5,
\end{equation}
the emitted power increases over the simulated working window. Averaging over twelve seeds, the early, middle, and late emitted-power means are
\begin{equation}
  0.0475,\qquad 0.0530,\qquad 0.0576.
\end{equation}
The middle-to-early acceleration ratio is therefore
\begin{equation}
  R_{\rm acc}=1.115.
\end{equation}
For the linear control, with the same channel count and coupling, the corresponding means are
\begin{equation}
  0.0854,\qquad 0.0783,\qquad 0.0732,
\end{equation}
and
\begin{equation}
  R_{\rm acc}=0.917.
\end{equation}

Thus the same class of Hamiltonian collision dynamics distinguishes the negative-heat-capacity profile from the linear profile: the convex case accelerates in the working window, while the linear control decelerates.

\begin{figure}[t]
  \centering
  \includegraphics[width=0.85\linewidth]{../step2_hamiltonian_polished.pdf}
  \caption{Polished Step 2 comparison. The convex core has negative microcanonical heat capacity over the working window; the linear control does not. Both use the same multi-mode collision-Hamiltonian architecture. The convex case shows an increasing emitted-power window, while the linear control decelerates. The radiation Renyi-2 entropy is computed from the reduced core state generated by the collision unitary.}
  \label{fig:step2-polished}
\end{figure}

\subsection{Parameter Scan}

A scan over curvature, emitted channel count, and coupling shows that the successful regime is restricted but not isolated. Acceleration occurs for sufficiently convex profiles, high outgoing channel count, and weak to moderate coupling. In particular, for curvature three and eight emitted labels the acceleration ratios are
\begin{equation}
  R_{\rm acc}=1.04,\ 1.17,\ 1.10,\ 0.40
\end{equation}
for
\begin{equation}
  g=0.3,\ 0.5,\ 0.8,\ 1.2,
\end{equation}
respectively. The linear controls do not accelerate. For curvature one and eight labels the corresponding ratios are
\begin{equation}
  R_{\rm acc}=0.94,\ 0.93,\ 0.75,\ 0.20.
\end{equation}
Strong coupling front-loads emission and washes out the negative-heat-capacity schedule. The correct regime is therefore weak-collision, many-channel, and finite-window.
